\documentclass[11pt]{article}
\usepackage[a4paper,margin=21mm]{geometry}
\usepackage{fontspec}
\setmainfont{Latin Modern Roman}
\newfontfamily\CJKfont{Noto Serif CJK SC}
\XeTeXlinebreaklocale "zh"
\XeTeXlinebreakskip=0pt plus 1pt
\usepackage{amsmath,amssymb,amsthm,amscd,graphicx,color}
\usepackage{xurl}
\usepackage[unicode,hidelinks]{hyperref}
\allowdisplaybreaks
\setlength\emergencystretch{3em}
\numberwithin{equation}{section}
\newtheorem{Theorem}{Theorem}[section]
\newtheorem*{Theorem*}{Theorem}
\newtheorem{Corollary}[Theorem]{Corollary}
\newtheorem{Lemma}[Theorem]{Lemma}
\newtheorem{Proposition}[Theorem]{Proposition}
 { \theoremstyle{definition}
\newtheorem{Definition}[Theorem]{Definition}
\newtheorem{Note}[Theorem]{Note}
\newtheorem{Example}[Theorem]{Example}
\newtheorem{Remark}[Theorem]{Remark} }

\def\N{\mathbb{N}}
\def\C{\mathbb{C}}% complex nember field
\def\R{\mathbb{R}}% real number field
\def\Q{\mathbb{Q}}% rational number field
\def\Z{\mathbb{Z}}% ring of integers
\def\A{\mathbb{A}}% the ring of adeles
\def\F{\mathbb{F}}% finite field
\def\H{\mathbb{H}}% Hamilton's quaternion field

%---------------Algebraic groups------------------------
\def\GL{\mathrm{GL}}
\def\PGL{\mathrm{PGL}}
\def\SL{\mathrm{SL}}
\def\PSL{\mathrm{PSL}}
%\def\O{\textnormal{O}}
\def\SO{\mathrm{SO}}
\def\GSO{\mathrm{GSO}}
\def\U{\textnormal{U}}
\def\SU{\textnormal{SU}}
\def\Sp{\mathrm{Sp}}
\def\GSp{\mathrm{GSp}}
\def\M{\mathrm{M}}% matrix space

\def\gl{\mathfrak{gl}}
\def\sl{\mathfrak{sl}}
\def\so{\mathfrak{so}}

%--------------symbols----------------------
\def\t{\,{}^t}% transpose

\def\d{\,\mathrm{d}}% measure

\def\idots{\reflectbox{$\ddots$}}% anti-diagonal dots

\def\Cc{C_c^{\infty}}% smooth functions

\def\ds{\displaystyle}% displaystyle

\def\bs{\backslash}% backslash

\def\t{\,{}^t}% transpose

\def\idots{\reflectbox{$\ddots$}}% anti-diagonal dots

\def\tbe{\textcolor{red}{\bf To Be Edited.}}

\DeclareMathOperator{\diag}{diag} % diagonal matrix

\DeclareMathOperator{\supp}{supp}% support

\DeclareMathOperator{\ord}{ord}% order

\DeclareMathOperator{\vol}{vol}% volume

\DeclareMathOperator{\val}{val}% valuation

\DeclareMathOperator{\mat}{Mat}% matrix space

\DeclareMathOperator{\rank}{rank}% rank

\DeclareMathOperator{\Gal}{Gal}% Galois group

\DeclareMathOperator{\Res}{Res}% residue / scalar restriction

\DeclareMathOperator{\Ind}{Ind}% induction

\DeclareMathOperator{\ind}{ind}% compact induction

\DeclareMathOperator{\Hom}{Hom}% space of homomorphisms

\DeclareMathOperator{\End}{End}% space of endomorphisms

\DeclareMathOperator{\Aut}{Aut}% group of automorphisms

\DeclareMathOperator{\Lie}{Lie}% Lie algebra

\DeclareMathOperator{\Inn}{Inn}% group of inner automorphisms

\DeclareMathOperator{\Out}{Out}% group of outer automorphisms

\DeclareMathOperator{\Cent}{Cent}% centralizer

\DeclareMathOperator{\Irr}{Irr}% set of irreducible representations

\DeclareMathOperator{\sgn}{sgn}% the sign character

\DeclareMathOperator{\re}{Re}% real part

\DeclareMathOperator{\im}{Im}% imaginary part


\DeclareMathOperator{\Ad}{Ad}% adjoint action
\DeclareMathOperator{\ad}{ad}% adjoint action
\DeclareMathOperator{\Kfin}{\text{$K$-fin}}% the space of K-finite vectors
\DeclareMathOperator{\Ker}{Ker}% kernel
\DeclareMathOperator{\Image}{Im}% image
\DeclareMathOperator{\triv}{triv}% the trivial representation
\DeclareMathOperator{\tr}{tr}% trace
\DeclareMathOperator{\id}{id}% identity operator
\DeclareMathOperator{\Sym}{Sym}% symmetrization
\DeclareMathOperator{\cpt}{cpt}% compact
\DeclareMathOperator{\semi}{ss}% semisimple

%------------fonts------------------------
%<<Roman>>
\def\a{\alpha}
\def\b{\beta}
\def\c{\gamma}
\def\d{\delta}
\def\e{\varepsilon}
\def\f{\varphi}
\def\g{\psi}
\def\h{\hbar}
%\def\i{\mbox{\raisebox{.5ex}{$\chi$}}}
\def\k{\kappa}
\def\l{\lambda}
\def\m{\mu}
\def\n{\nu}
%\def\o{\omega}
\def\ro{\rho}
\def\s{\sigma}
\def\t{\tau}
\def\x{\xi}
\def\y{\eta}
\def\z{\zeta}
\def\th{\theta}
\def\pa{\partial}
\newcommand{\jap}[1]{\langle #1 \rangle}
%<<mathfrak>>
\newcommand{\fa}{\mathfrak{a}}
\newcommand{\fb}{\mathfrak{b}}
\newcommand{\fc}{\mathfrak{c}}
\newcommand{\fe}{\mathfrak{e}}
\newcommand{\ff}{\mathfrak{f}}
\newcommand{\fg}{\mathfrak{g}}
\newcommand{\fh}{\mathfrak{h}}
\newcommand{\fii}{\mathfrak{i}}
\newcommand{\fj}{\mathfrak{j}}
\newcommand{\fk}{\mathfrak{k}}
\newcommand{\fl}{\mathfrak{l}}
\newcommand{\fm}{\mathfrak{m}}
\newcommand{\fn}{\mathfrak{n}}
\newcommand{\fo}{\mathfrak{o}}
\newcommand{\fp}{\mathfrak{p}}
\newcommand{\fq}{\mathfrak{q}}
\newcommand{\fr}{\mathfrak{r}}
\newcommand{\fs}{\mathfrak{s}}
\newcommand{\fS}{\mathfrak{S}}
\newcommand{\ft}{\mathfrak{t}}
\newcommand{\fu}{\mathfrak{u}}
\newcommand{\fv}{\mathfrak{v}}
\newcommand{\fx}{\mathfrak{x}}
\newcommand{\fy}{\mathfrak{y}}
\newcommand{\fz}{\mathfrak{z}}
\newcommand{\fI}{\mathfrak{I}}
\newcommand{\fJ}{\mathfrak{J}}
\newcommand{\fK}{\mathfrak{K}}
\newcommand{\fL}{\mathfrak{L}}
\newcommand{\fP}{\mathfrak{P}}
\newcommand{\fT}{\mathfrak{T}}
\newcommand{\fW}{\mathfrak{W}}
\newcommand{\fX}{\mathfrak{X}}

%<< mathbb>>
%\newcommand{\bA}{\mathbb{A}}
\newcommand{\bB}{\mathbb{B}}
%\newcommand{\bC}{\mathbb{C}}
\newcommand{\bD}{\mathbb{D}}
\newcommand{\bE}{\mathbb{E}}
\newcommand{\bF}{\mathbb{F}}
\newcommand{\bG}{\mathbb{G}}
%\newcommand{\bH}{\mathbb{H}}
\newcommand{\bI}{\mathbb{I}}
\newcommand{\bJ}{\mathbb{J}}
\newcommand{\bK}{\mathbb{K}}
\newcommand{\bL}{\mathbb{L}}
\newcommand{\bM}{\mathbb{M}}
%\newcommand{\bN}{\mathbb{N}}
\newcommand{\bO}{\mathbb{O}}
\newcommand{\bP}{\mathbb{P}}
%\newcommand{\bQ}{\mathbb{Q}}
%\newcommand{\bR}{\mathbb{R}}
\newcommand{\bS}{\mathbb{S}}
\newcommand{\bT}{\mathbb{T}}
\newcommand{\bU}{\mathbb{U}}
\newcommand{\bV}{\mathbb{V}}
\newcommand{\bW}{\mathbb{W}}
\newcommand{\bX}{\mathbb{X}}
\newcommand{\bY}{\mathbb{Y}}
%\newcommand{\bZ}{\mathbb{Z}}

%<<mathcal>>
\newcommand{\cA}{\mathcal{A}}
\newcommand{\cB}{\mathcal{B}}
\newcommand{\cC}{\mathcal{C}}
\newcommand{\cD}{\mathcal{D}}
\newcommand{\cE}{\mathcal{E}}
\newcommand{\cF}{\mathcal{F}}
\newcommand{\cG}{\mathcal{G}}
\newcommand{\cH}{\mathcal{H}}
\newcommand{\cI}{\mathcal{I}}
\newcommand{\cJ}{\mathcal{J}}
\newcommand{\cK}{\mathcal{K}}
\newcommand{\cL}{\mathcal{L}}
\newcommand{\cM}{\mathcal{M}}
\newcommand{\cN}{\mathcal{N}}
\newcommand{\cO}{\mathcal{O}}
\newcommand{\cP}{\mathcal{P}}
\newcommand{\cQ}{\mathcal{Q}}
\newcommand{\cR}{\mathcal{R}}
\newcommand{\cS}{\mathcal{S}}
\newcommand{\cT}{\mathcal{T}}
\newcommand{\cU}{\mathcal{U}}
\newcommand{\cV}{\mathcal{V}}
\newcommand{\cW}{\mathcal{W}}
\newcommand{\cX}{\mathcal{X}}
\newcommand{\cY}{\mathcal{Y}}
\newcommand{\cZ}{\mathcal{Z}}


\newcommand{\sI}{\mathscr{I}}

\newcommand{\bfe}{\mathbf{e}}
\newcommand{\bfh}{\mathbf{h}}
\newcommand{\bfi}{\mathbf{i}}
\newcommand{\bfj}{\mathbf{j}}
\newcommand{\bfk}{\mathbf{k}}
\newcommand{\bfI}{\mathbf{I}}
\newcommand{\bfJ}{\mathbf{J}}
\newcommand{\bfK}{\mathbf{K}}


\begin{document}
\begin{center}\Large Uniform symbol expansion of Gegenbauer polynomials\end{center}
\noindent\textbf{Stable ID:} 554\quad\textbf{Author:} Kouichi Taira; Hiroyoshi Tamori\par
\noindent\textbf{Work:} Strichartz Estimates for the (k,a)-Generalized Laguerre Operators\par
\noindent\textbf{Source:} \url{https://sigma-journal.com/2025/014/}\par
\noindent\textbf{Version:} Official SIGMA 21 (2025) PDF and native TeX; acquired 2026-09-30; exact bytes pinned by SHA256\par
\noindent\textbf{Rights:} CC BY 4.0. Authors retain copyright. \url{https://creativecommons.org/licenses/by/4.0/}. Reuse and typesetting adaptation; original language and mathematical content preserved.\par
\noindent\textbf{Difficulty (prior selection preserved):} {\CJKfont H2: The polynomial degree must be extended to a smooth real parameter while preserving endpoint-uniform symbol derivatives. Branch selection, singular-endpoint amplitudes, and the degeneration of the phase derivative as the angle approaches 0 or pi require a parameter-scaled construction and rapidly decreasing remainder, rather than substituting a standard asymptotic formula.}\par
\medskip\noindent\textbf{Editorial boundary note (not author text):} Complete original Proposition 3.3 on uniform Gegenbauer symbol expansion, including the full Beta-factor derivative estimates, branch-aware integral decomposition, endpoint/interior remainder regimes, and the actual symbol-estimate lemma and proof. The uniform singular-amplitude nonstationary argument and Leibniz bound are retained. The author explicitly states that this extends earlier normalized estimates in parameter range and includes the normalization factor. It is a family-wide special-function asymptotic result, distinct from the moving-critical-point theorem; no separate count for the Beta-factor technical proposition.\par
\medskip\hrule\medskip\noindent\textbf{Author text begins below.}\par
\setcounter{section}{1}\setcounter{subsection}{4}
% Original source lines 506--511
\subsection{Notation}
Let
\begin{align*}
\mathbb{N}=\{0,1,2,\hdots\},\qquad \mathbb{N}^*=\{1,2,3,\hdots\}, \qquad \R_{+}=(0,\infty),\qquad \R_-=(-\infty,0).
\end{align*}
For a parameter $i\in I$ and $A_i,B_i$, we write $A_i\lesssim B_i$ for $i\in I$ if there is a constant $C>0$ independent of $i\in I$ such that $A_i\leq CB_i$. We denote $A_i\sim B_i$ if both $A_i\lesssim B_i$ and $B_i\lesssim A_i$ hold for $i\in I$.

\setcounter{section}{3}\setcounter{subsection}{1}\setcounter{Theorem}{2}
% Original source lines 743--777
\subsection{Asymptotic expansion of Gegenbauer polynomials}

The following proposition is a refinement of \cite[Lemma 10.2]{IJ} and its proof is given in Appendix~\ref{appasymgegensubsec}.

\begin{Proposition}\label{Gegenasymp}
Let $\n\geq 0$.
There are functions $g_{+}(m,\f), g_{-}(m,\f)$ which are smooth respect to $m\in [1,\infty)$ and a function $r(m,\f)$ defined for $m\in\mathbb{N}^*$ such that
\begin{align*}
\n^{-1}C_m^{\n}(\cos\f)=\sum_{\pm}g_{\n,\pm}(m,\f){\rm e}^{\pm {\rm i}m\f}+r(m,\f)
\end{align*}
for $m\in \mathbb{N}^*$ $($we interpret the left-hand side for $\n=0$ as \smash{$\lim_{\n\searrow 0}\n^{-1}C_m^{\n}(\cos \f))$}, and
\begin{align*}
&|\pa_m^{\a}g_{\n,\pm}(m,\f)|\leq
C_{\a,\n}m^{2\n-1-\a}(1+m\sin\f)^{-\n},\qquad m\geq 1,\\
&|r(m,\f)|\leq C_{N,\n}m^{-N},\qquad m\in\mathbb{N}^*
\end{align*}
for each $N>0$, $\a\in\mathbb{N}$, and $\f\in [0,\pi]$ with constants $C_{\a,\n}, C_{N,\n}>0$.
\end{Proposition}

\begin{Remark}\quad
\begin{itemize}\itemsep=0pt
\item[(1)] In particular, for each $\e>0$, we have $|\pa_m^{\a}g_{\n,\pm}(m,\f)|\leq C_{\a,\n,\e}m^{\n-1-\a}$ uniformly in $\f\in [\e,\pi-\e]$.
\item[(2)] We do not impose that $g_{\n,\pm}$ are continuous with respect to $\f\in [0,\pi]$. In this paper, we just need uniform estimates of $g_{\n,\pm}$ in $\f\in [0,\pi]$.
\item[(3)]
In \cite[Lemma 10.2]{IJ}, a similar estimate for $\n^{-1}C_m^{\n}(\cos\f)$ divided by $C_m^{\n}(1)$ is given.
Here we also deal with the estimate for $C_m^{\n}(1)$.
Moreover, the ranges of the parameters are extended compared to there.
\end{itemize}
\end{Remark}

In particular, we have a uniform bound: For $\n\geq 0$, there exists $C>0$ such that
\begin{align}\label{Gegenunif}
\big|\n^{-1}C_m^{\n}(\cos \f)\big|\leq Cm^{2\n-1}\qquad m\in \mathbb{N}^*,\quad \f\in[0,\pi],
\end{align}
where the left-hand side is interpreted as $\big|\lim_{\n\searrow 0}\n^{-1}C_m^{\n}(\cos \f)\big|$ when $\n=0$.

\appendix\setcounter{section}{0}
% Original source lines 1539--1599
\section{Asymptotic behavior of special functions}

\subsection{Leibniz's rule}

We frequently use the following formula, which directly follows from Leibniz's rule.

\begin{Lemma}\label{Leibnizrule}
Let $r$ be a smooth function and $N\in\mathbb{N}\setminus\{0\}$. Then there exist smooth functions~$b_{jN}$ such that
\begin{gather*}
\bigl(\pa_x\circ r(x)^{-1}\bigr)^N=\frac{1}{r(x)^N}\!\sum_{j=0}^Nb_{jN}(x)\pa_{x}^j\qquad \text{with}\quad |b_{jN}(x)|\lesssim \sum_{k=1}^{N-j}\sum_{\substack{\ell_1+\cdots+\ell_k=N-j,\\ \ell_1,\hdots,\ell_k\geq 1}}\prod_{i=1}^k\!\frac{\big|r^{(\ell_i)}(x)\big|}{|r(x)|}.
\end{gather*}
\end{Lemma}

\begin{Remark}
More precisely, we have
\begin{align*}
\bigl(\pa_x\circ r(x)^{-1}\bigr)^N=\frac{1}{r(x)^N}\pa_x^N+\frac{1}{r(x)^N}\sum_{j=0}^{N-1}\Bigg(\sum_{k=1}^{N-j}\sum_{\substack{\ell_1+\cdots+\ell_k=N-j,\\ \ell_1,\hdots,\ell_k\geq 1}}C_{\ell_1\ell_2\hdots \ell_k}^{jN}\prod_{i=1}^{k}\frac{r^{(\ell_i)}(x)}{r(x)} \Bigg)\pa_{x}^j
\end{align*}
with constants \smash{$C_{\ell_1\ell_2\hdots \ell_m}^{jN}>0$}.
\end{Remark}

\subsection{Non-stationary phase theorem with a singular amplitude}

The following lemma is more or less well known and an easy consequence of integration by parts. We give a proof for completeness of the paper.

\begin{Lemma}\label{singularasym}
Let $\m>-1$, $\c\in C^{\infty}((0,\infty))$ which is supported in $x\leq 10$ such that for $\a\in\mathbb{N}$, we have $|\pa_x^{\a}\c(x)|\lesssim |x|^{\m-\a}$ for $x\in (0,\infty)$. Let $f\in C^{\infty}(\R)$ satisfy $\im f(x)\geq 0$, $f'(x)\neq 0$ for~${x\in \supp \c}$ and $f(0)\in \R$. We define
\begin{align*}
b_{\m}(\l):={\rm e}^{-{\rm i}\l f(0)}\int_{0}^{\infty}\c(x){\rm e}^{{\rm i}\l f(x)}{\rm d}x\qquad \text{for}\quad \l\gtrsim 1.
\end{align*}
Then for each $\a\in\mathbb{N}$
\begin{align}\label{singbmues}
|\partial_{\l}^{\a}b_{\m}(\l)|\lesssim \l^{-\m-1-\a},\qquad \l\gtrsim 1.
\end{align}
In particular, $\big|\int_{0}^{\infty}\c(x){\rm e}^{{\rm i}\l f(x)}{\rm d}x\big|\lesssim \l^{-\m-1}$ for $\l\gtrsim 1$.
\end{Lemma}

\begin{proof}
We only need to prove these estimates for sufficiently large $\l$.

First, we prove \eqref{singbmues} for $\a=0$. Let $\chi\in C_c^{\infty}(\R)$ such that $\chi(x)=1$ for $|x|\leq 1/2$ and~${\chi(x)=0}$ for $|x|\geq 1$. Define $\chi_{\l}(x)=\chi(\l x)$ and $\overline{\chi_\l}(x)=1-\chi_{\l}(x)$. Then we have
\begin{align}\label{singbmues1}
\left|\int_{0}^{\infty}\c(x)\chi_{\l}(x){\rm e}^{{\rm i}\l f(x)}{\rm d}x\right|\lesssim \int_{0}^{1/\l}x^{\m}{\rm d}x\lesssim {\l}^{-\m-1}.
\end{align}
On the other hand, the integration by parts yields
\begin{align*}
\left|\int_{0}^{\infty}\c(x)\overline{\chi_{\l}}(x){\rm e}^{{\rm i}\l f(x)}{\rm d}x\right|=\l^{-N}\left|\int_{0}^{\infty}L^N(\c(x)\overline{\chi_{\l}}(x)){\rm e}^{{\rm i}\l f(x)}{\rm d}x\right|,
\end{align*}
where we define $L=\pa_x\circ ({\rm i}f'(x))^{-1}$. By using Leibniz's rule, Lemma \ref{Leibnizrule} and the assumption $f'\neq 0$ on $\supp \c$, we obtain $|L^{N}(\c(x)\overline{\chi_\l}(x))|\lesssim |x|^{\m-N}$. Hence, for $N>\m+1$,
\begin{align}\label{singbmues2}
\left|\int_{0}^{\infty}\c(x)\overline{\chi_{\l}}(x){\rm e}^{{\rm i}\l f(x)}{\rm d}x\right|\lesssim \l^{-N}\int_{\frac{1}{2\l}}^{10}x^{\m-N}{\rm d}x \lesssim \l^{-\m-1}.
\end{align}
The inequalities \eqref{singbmues1} and \eqref{singbmues2} imply \eqref{singbmues} for $\a=0$.

To deal with the case $\a\geq 1$, we write
\begin{align*}
\pa_{\l}^{\a}b_{\m}(\l)={\rm i}^{\a}{\rm e}^{-{\rm i}\l f(0)}\int_0^{\infty}(f(x)-f(0))^{\a}\c(x){\rm e}^{{\rm i}\l f(x)}{\rm d}x.
\end{align*}
By Taylor's theorem, we have $\pa_x^{N}(\c(x)(f(x)-f(0))^{\a})=O\bigl(|x|^{\m+\a-N}\bigr)$ for $x\in \supp \c$. Moreover, $\big|{\rm e}^{-{\rm i}\l f(0)}\big|=1$ due to the assumption $f(0)\in \R$.
Hence the same argument as the case $\a=0$ gives \eqref{singbmues} for $\a\geq 1$.
\end{proof}

\setcounter{subsection}{4}\setcounter{Theorem}{9}
% Original source lines 1901--2073
\subsection{Asymptotics of the Beta function}\label{subsecasymbeta}

For $m\in [1,\infty)$ and $\n>0$, we define
\begin{align}\label{Fnu1def}
F_{\n,1}(m):=\n^{-1}\frac{\Gamma\bigl(\n+\frac{1}{2}\bigr)}{\sqrt{\pi}\Gamma(\n)}\cdot \frac{\Gamma(m+2\n)}{\Gamma(m+1)\Gamma(2\n)}.
\end{align}
By \cite[Fact 4.8, equation~(4.30)]{BKO}, we have $C_m^{\n}(1)=\frac{\Gamma(m+2\n)}{\Gamma(m+1)\Gamma(2\n)}$ and hence
\begin{align}\label{Fnu1Gegen}
F_{\n,1}(m)=\n^{-1}\frac{\Gamma\bigl(\n+\frac{1}{2}\bigr)C_m^{\n}(1)}{\sqrt{\pi}\Gamma(\n)}\qquad \text{for}\quad m\in\mathbb{N}.
\end{align}
We denote the Beta function by $B(x,y)$. Then it is known that $B(x,y)=\frac{\Gamma(x)\Gamma(y)}{\Gamma(x+y)}$. We may write
\begin{align}\label{Fnu1beta}
F_{\n,1}(m)=\n^{-1}\frac{\Gamma\bigl(\n+\frac{1}{2}\bigr)}{\sqrt{\pi}\Gamma(\n)}\cdot \frac{1}{mB(m,2\n)}.
\end{align}

\begin{Proposition}\label{Fnuestimate}
Let $\a\geq 0$ be an integer and $\n>0$. Then there exists $C_{\n,\a}>0$ such that
\begin{align*}
|\pa_m^{\a}B(m,2\n)|\leq C_{\n,\a}m^{-2\n-\a},\qquad |\pa_m^{\a}F_{\n,1}(m)|\leq C_{\n,\a}m^{2\n-1-\a}
\end{align*}
for all $m\in [1,\infty)$.
\end{Proposition}

\begin{proof}First, we note that it suffices to prove these estimates for sufficiently large $m$. As is well known (essentially due to Stirling's formula), we have $B(m,2\n)\sim \Gamma(2\n)m^{-2\n}$ as $m\to \infty$.
This implies that the estimates for $F_{\n,1}(m)$ follow from the estimates for $B(m,2\n)$ and \eqref{Fnu1beta}. The case $\a=0$ follows from $B(m,2\n)\sim \Gamma(2\n)m^{-2\n}$ as $m\to \infty$. Hence, in the following, we consider the estimates for $\a\geq 1$.

Let $\g\in C^{\infty}(\R;[0,1])$ satisfy $\g(t)=1$ for $t\geq 1/2$ and $\g(t)=0$ for $t\leq \frac{1}{4}$.
Since $B(m,2\n)=\int_0^1t^{m-1}(1-t)^{2\n-1}{\rm d}t$, we have
\begin{align*}
\pa_m^{\a}B(m,2\n)=\int_0^1t^{m-1}g_{\a,1}(t){\rm d}t+\int_0^1t^{m-1}g_{\a,2}(t){\rm d}t,
\end{align*}
where we set $g_{\a,1}(t)=(\log t)^{\a}(1-t)^{2\n-1}\g(t)$ and $g_{\a,2}(t)=(\log t)^{\a}(1-t)^{2\n-1}(1-\g(t))$.

First, we deal with the second term. Since $\supp g_{\a,2}\subset \{t\leq 1/2\}$ and since $|(\log t)^{\a}|\lesssim |t|^{-1}$ for $0<t\leq 1/2$, we have
\begin{align*}
\left|\int_0^1t^{m-1}g_{\a,2}(t){\rm d}t\right|\lesssim \int_0^{\frac{1}{2}}t^{m-2}{\rm d}t\lesssim 2^{-m}\lesssim m^{-2\n-\a}
\end{align*}
for sufficiently large $m\geq 1$.

Next, we consider the first term. We note
\begin{align*}
\big|\pa_{t}^{\b}g_{\a,1}(t)\big|\leq C_{\a}|1-t|^{2\n-1+\a-\b}\qquad\text{for}\quad t\in [0,1].
\end{align*}
Now we write
\begin{align*}
t^{m-1}={\rm e}^{(m-1)(\log t)}={\rm e}^{{\rm i}(m-1)\tilde{f}(t)},\qquad \tilde{f}(t)=-{\rm i}\log t.
\end{align*}
Then $\tilde{f}'(t)=-{\rm i}/t\neq 0$ and $\im \tilde{f}(t)\geq 0$ hold for $t\in \supp \g\cap (-\infty,1]$.
Now it follows from Lemma \ref{singularasym} with $\l=m$, $\m=2\n-1$ and the change of variable $x=1-t$ that
\begin{align*}
\left|\int_0^1t^{m-1}g_{\a,1}(t){\rm d}t\right|\lesssim m^{-2\n-\a}
\end{align*}
holds for sufficiently large $m\geq 1$. This completes the proof.
\end{proof}

\subsection{Asymptotics of the Gegenbauer polynomials}\label{appasymgegensubsec}

In this appendix, we prove Proposition \ref{Gegenasymp} using non-stationary phase theorem. The case~${\n=0}$ is easy to deal with (see the proof of Proposition \ref{Gegenasymp} below) and hence we focus on the case~${\n>0}$.

By \cite[equation~(4.30)]{BKO} (see also \cite[Section 10]{IJ}) and \eqref{Fnu1Gegen},
\begin{align}\label{Gegenintrep}
\n^{-1}C_m^{\n}(\cos\f)=&F_{\n,1}(m)\int_{-1}^{1}(\cos\f+{\rm i}(\sin\f)u)^m\bigl(1-u^2\bigr)^{\n-1}{\rm d}u
\end{align}
for $\n>0$, $m\in\mathbb{N}^*$ and $\f\in [0,\pi]$, where $F_{\n,1}$ is defined in \eqref{Fnu1def}. To get an asymptotics of~$C_m^{\n}(\cos\f)$, we extend $\mathbb{N}^*\ni m\mapsto C_m^{\n}(\cos\f)$ to a function on $m\in [1,\infty)$ (up to a negligible term). However, since the integrant $(\cos\f+{\rm i}(\sin\f)u)^m$ is not a single-valued function, we have to do it a bit carefully.

First, we deal with the case $\f\in \bigl(\frac{\pi}{4}, \frac{3}{4}\pi\bigr)$.
We consider two branches of logarithmic smooth functions $\log^+z\colon\mathbb{C}\setminus \{{\rm i}y\mid y\leq 0\}\to \mathbb{C}$ and $\log^-z\colon\mathbb{C}\setminus \{{\rm i}y\mid y\geq 0\}\to \mathbb{C}$ such that $\log^+(\cos \f+{\rm i}\sin \f)={\rm i}\f$ and $\log^-(\cos\f-{\rm i}\sin\f))=-{\rm i}\f$ for all $\f\in [0,\pi]$. We define smooth functions~$f_{\pm}$~by
\begin{align*}
f_{\pm}(\f,u)=-{\rm i}\log^{\pm}(\cos\f+{\rm i}(\sin\f)u)\mp \f
\end{align*}
for $(\f,u)$ such that $\cos\f+{\rm i}(\sin\f)u$ belongs to the domain of $\log^{\pm}$. We note that
$([0,\pi]\times \R_{\pm})\cup \bigl(\bigl(v[0,\frac{\pi}{4}\bigr]\cup \bigl[\frac{3}{4}\pi,\pi\bigr]\bigr)\times \bigl[-\frac{1}{2},\frac{1}{2}\bigr] \bigr)$ is included in the domain of $f_{\pm}$.
Then ${\rm e}^{{\rm i}mf_{\pm}(\f,u)}=(\cos\f+{\rm i}(\sin\f)u)^m{\rm e}^{\mp {\rm i}m\f}$ for all $m\in\mathbb{N}^*$ and $\pm u>0$.
Now let $\chi_{\pm},\chi_0\in C^{\infty}(\R;[0,1])$ such that $\chi_++\chi_0+\chi_-=1$ on $\R$, $\supp\chi_0\subset \bigl[-\frac{1}{2},\frac{1}{2}\bigr]$, $\supp \chi_+\subset \bigl[\frac{1}{4},\infty\bigr)$ and $\supp\chi_-\subset \bigl(-\infty,-\frac{1}{4}\bigr]$. We define
\begin{align*}
&H_{\n,\pm}(m,\f):=\int_{-1}^1{\rm e}^{{\rm i}mf_{\pm}(\f,u)}\chi_{\pm}(u)\bigl(1-u^2\bigr)^{\n-1}{\rm d}u,\\
&E_{\n}(m,\f):=\int_{-1}^1(\cos \f+{\rm i}(\sin\f)u)^m\chi_{0}(u)\bigl(1-u^2\bigr)^{\n-1}{\rm d}u,
\end{align*}
where $H_{\n,\pm}(m,\f)$ is defined for all $m\in [1,\infty)$ although $E_{\n}(m,\f)$ is defined only for $m\in \mathbb{N}^*$ whenever $\f\in \bigl(\frac{\pi}{4}, \frac{3}{4}\pi\bigr)$ (note that $\cos\f+{\rm i}(\sin \f)u$ may be zero for $u\in \supp \chi_0$). By these construction and \eqref{Gegenintrep}, we have
\begin{align}\label{Gegendecom}
\n^{-1}C_m^{\n}(\cos\f)=&F_{\n,1}(m)\bigl({\rm e}^{{\rm i}m\f}H_{\n,+}(m,\f)+{\rm e}^{-{\rm i}m\f}H_{\n,-}(m,\f)+E_{\n}(m,\f) \bigr)
\end{align}
for $m\in\mathbb{N}^*$ and $\f\in \bigl(\frac{\pi}{4}, \frac{3}{4}\pi\bigr)$.

Next, we consider the case $\f\in \bigl[0,\frac{\pi}{4}\bigr]\cup \bigl[\frac{3}{4}\pi,\pi\bigr]$. We define
\begin{gather*}
E_{\n}'(m,\f):=\int_{-1}^1{\rm e}^{{\rm i}mf_{+}(\f,u)}\chi_{0}(u)\bigl(1-u^2\bigr)^{\n-1}{\rm d}u,\qquad m\in [1,\infty),\\ \f\in \left[0,\frac{\pi}{4}\right]\cup \left[\frac{3}{4}\pi,\pi\right],
\end{gather*}
where we note that $\cos \f+{\rm i}(\sin\f)u\in \mathbb{C}\setminus \{{\rm i}y\mid y\leq 0\}$ whenever $\f\in \bigl[0,\frac{\pi}{4}\bigr]\cup \bigl[\frac{3}{4}\pi,\pi\bigr]$ and $u\in \supp \chi_0$. Then
\begin{align}\label{Gegendecom2}
\n^{-1}C_m^{\n}(\cos\f)=&F_{\n,1}(m)\bigl({\rm e}^{{\rm i}m\f}H_{\n,+}(m,\f)+{\rm e}^{-{\rm i}m\f}H_{\n,-}(m,\f)+{\rm e}^{{\rm i}m\f}E_{\n}'(m,\f) \bigr)
\end{align}
holds for $m\in\mathbb{N}^*$ and $\f\in \bigl[0,\frac{\pi}{4}\bigr]\cup \bigl[\frac{3}{4}\pi,\pi\bigr]$.

Now we show that $H_{\n,\pm}$, $E_{\n}$ and $E_{\n}'$ satisfy symbol-type estimates. The basic idea of the proof for $H_{\n,\pm}$ and $E_{\n}$ is to use the non-stationary phase theorem with a singular amplitude (see Lemma~\ref{singularasym}). To do this, we observe that our phase function $f_{\pm}(\f,u)$ satisfies $|f_{\pm}'(\f,u)|\gtrsim |\sin \f|$ on~${u\in \supp \chi_{\pm}}$ (see the proof below). Although $\pa_u f_{\pm}(\f,u)=0$ when $\f=0,\pi$, we consider a~large parameter $\l=m\sin \f$ rather than $m$ and can apply Lemma \ref{singularasym}.

\begin{Lemma}\label{Gegenlemma}
Let $\a\geq 0$ be an integer and $\n,N>0$.
\begin{itemize}\itemsep=0pt
\item[$(i)$] There exists $C_{N,\n}>0$ such that $|E_{\n}(m,\f)|\leq C_{N,\n}m^{-N}$ for $m\in\mathbb{N}^*$ and $\f\in \bigl(\frac{\pi}{4}, \frac{3}{4}\pi\bigr)$.

\item[$(ii)$] There exists $C_{\a,N,\n}>0$ such that $|\pa_{m}^{\a}E_{\n}'(m,\f)|\leq C_{\a,N,\n}m^{-\a}(1+m\sin\f)^{-N}$ for $m\in[1,\infty)$ and $\f\in \bigl[0,\frac{\pi}{4}\bigr]\cup \bigl[\frac{3}{4}\pi,\pi\bigr]$.

\item[$(iii)$] There exists $C_{\a,\n}>0$ such that $|\pa_{m}^{\a}H_{\n,\pm}(m,\f)|\leq C_{\a,\n}m^{-\a}(1+m\sin\f)^{-\n}$ for $m\in[1,\infty)$ and $\f\in [0,\pi]$.
\end{itemize}
\end{Lemma}

\begin{proof}
 $(i)$ Since \smash{$|\cos\f+{\rm i}(\sin \f)u|\leq \sqrt{\frac{5}{8}}$} for $u\in \supp \chi_0$ and $\f\in \bigl(\frac{\pi}{4}, \frac{3}{4}\pi\bigr)$, we have
 \[
|\cos\f+{\rm i}(\sin \f)u|^m\leq \left(\frac{5}{8}\right)^{\frac{m}{2}}\lesssim m^{-N}
\]
 for $m\in\mathbb{N}^*$. The estimate for $E_{\n}$ directly follows from this inequality.

 $(ii)$ Setting $\g(u):=\chi_0(u)\bigl(1-u^2\bigr)^{\n-1}$, we write $E_{\n}'(m,\f)=\int_{\R}{\rm e}^{{\rm i}mf_{+}(\f,u)}\g(u){\rm d}u$.

First, we consider the case $\a=0$ by using the integration by parts.
Since $\log^+(\cos \f+{\rm i}\sin \f)={\rm i}\f$, we have
\begin{align}\label{f_+rep}
f_+(\f,u)={\rm i}\int_u^1\frac{{\rm d}}{{\rm d}r}\log^+(\cos \f+{\rm i}(\sin \f) r){\rm d}r=-\int_u^1\frac{\sin\f(\cos \f-{\rm i}(\sin \f) r)}{\cos^2 \f+r^2\sin^2 \f }{\rm d}r
\end{align}
and hence \smash{$\im f_+(\f,u)=\int_u^1\frac{(\sin^2 \f) r}{\cos^2 \f+r^2\sin^2 \f }{\rm d}r$}. This implies $\im f_+(\f,u)\geq 0$ for $u\in [-1,1]$. In fact, this is trivial for $u\geq 0$. For $u\leq 0$, we write
\begin{align*}
\int_u^1\frac{\bigl(\sin^2 \f\bigr) r}{\cos^2 \f+r^2\sin^2 \f }{\rm d}r={}&\int_0^1\frac{\bigl(\sin^2 \f\bigr) r}{\cos^2 \f+r^2\sin^2 \f }{\rm d}r+\int_u^0\frac{\bigl(\sin^2 \f\bigr) r}{\cos^2 \f+r^2\sin^2 \f }{\rm d}r\\
={}&\int_{-u}^1\frac{\bigl(\sin^2 \f\bigr) r}{\cos^2 \f+r^2\sin^2 \f }{\rm d}r,
\end{align*}
 which is non-negative.

Now we set $L=\pa_u\circ ({\rm i}\pa_uf_+(\f,u))^{-1}$. By integrating by parts, we have
\[
|E_{\n}'(m,\f)|=m^{-N}\left|\int_{\R}{\rm e}^{{\rm i}mf_{+}(\f,u)} L^N\g(u){\rm d}u\right|\leq m^{-N}\int_{\R}\big| L^N\g(u)\big|{\rm d}u
\]
 by virtue of $\im f_+(\f,u)\geq 0$. Since $\pa_u^{\a}f_+(\f,u)=(-1)^{\a-1}(\a-1)!(\sin\f)^{\a}(\cos \f+{\rm i}u\sin \f )^{-\a}$, for each $\a\geq 2$, we obtain $|\pa_uf_+(\f,u)|^{-1}\lesssim |\sin \f|^{-1}$ and $|\pa_u^{\a}f_+(\f,u)||\pa_uf_+(\f,u)|^{-1} \lesssim 1$ for $u\in [-1,1]$ and $\f\in \bigl[0,\frac{\pi}{4}\bigr]\cup \bigl[\frac{3}{4}\pi,\pi\bigr]$. By Lemma \ref{Leibnizrule}, we conclude $\big| L^N\g(u)\big|\lesssim |\sin\f|^{-N}$ and hence $|E_{\n}'(m,\f)|\lesssim m^{-N}|\sin\f|^{-N}$ for $\f\in \bigl[0,\frac{\pi}{4}\bigr]\cup \bigl[\frac{3}{4}\pi,\pi\bigr]$ (we note that $\g$ is compactly supported). On the other hand, the inequality $\im f_+(\f,u)\geq 0$ also implies $|E_{\n}'(m,\f)|\lesssim 1$. Combining these estimates, we have proved (ii) for $\a=0$.

Next, we consider the case $\a\geq 1$. We note
\begin{align*}
\pa_m^{\a}E'_{\n}(m,\f)={\rm i}^{\a}\int_{\R}{\rm e}^{{\rm i}mf_{+}(\f,u)}f_{+}(\f,u)^{\a}\g(u)du.
\end{align*}
As in the case of $\a=0$, we can deduce $\big|L^N(f_{+}(\f,u)^{\a}\g(u))\big|\lesssim |{\sin\f}|^{\a-N}$, which leads to the part (ii) for $\a\geq 1$ since $f_{+}(0,u)=f_{+}(\pi,u)=0$.

 $(iii)$ We deal with the case $+$ only. There exists $c>0$ such that $|{\cos\f+{\rm i}u\sin \f}|\geq c$ for~${u\in \supp\chi_+}$ and $\f\in [0,\pi]$.
 Since \smash{$\im f_+(\f,u)=\int_u^1\frac{(\sin^2 \f) r}{\cos^2 \f+r^2\sin^2 \f }{\rm d}r$}, which is proved above, we have $\im f_+(\f,u)\geq 0$ for $u\in \supp \chi_+$ and $\f\in [0,\pi]$. Moreover, we have $f_+(\f,1)=0$.

Now we prove $(iii)$ for $\a=0$. Set $m'=m\sin \f$ and $g(\f,u)=(\sin \f)^{-1}f_+(\f,u)$. Then we write \smash{$H_{\n,+}(m,\f)=\int_{-1}^1{\rm e}^{{\rm i}m'g(\f,u)}\chi_{+}(u)\bigl(1-u^2\bigr)^{\n-1}{\rm d}u$}. Since
\[
\pa_u^{\a}g(\f,u)=(-1)^{\a-1}(\a-1)!(\sin\f)^{\a-1}(\cos \f+{\rm i}u\sin \f )^{-\a}
\]
 for $\a\geq 1$ and since $|{\cos \f+{\rm i}u\sin \f}|\in \bigl[\frac{1}{\sqrt{2}},1\bigr]$ for $\f\in \bigl[0,\frac{\pi}{4}\bigr]\cup \bigl[\frac{3}{4}\pi,\pi\bigr]$, we have
\begin{align*}
|\pa_ug(\f,u)|\geq c,\qquad |\pa_u^{\a}g(\f,u)| |\pa_ug(\f,u)|^{-1}\leq C_{\a}'
\end{align*}
with a constant $c>0$ and $C_{\a}'>0$. By Lemma \ref{singularasym} with $\l=m\sin \f$ and the change of variable~${x=1-u}$, we obtain $|H_{\n,+}(m,\f)|\lesssim (m')^{-\n}$ for $m'=m\sin\f\geq 1$. On the other hand, it follows from $\im f_+(\f,u)\geq 0$ that $|H_{\n,+}(m,\f)|\lesssim 1$. Combining these estimates, we obtain~${|H_{\n,+}(m,\f)|\lesssim (1+m')^{-\n}=(1+m\sin\f)^{-\n}}$ for $m\sin \f\geq 1$. For $m\sin \f\leq 1$, this estimate is easy to prove.

Finally, we consider the case $\a\geq 1$. We observe that
\begin{align*}
\pa_m^{\a}H_{\n,+}(m,\f)=i^{\a}\int_{\R}{\rm e}^{{\rm i}m'g(\f,u)}f_{+}(\f,u)^{\a}\bigl(1-u^2\bigr)^{\n-1}\chi_+(u){\rm d}u
\end{align*}
and \smash{$\big|\pa_u^{\b}f_{+}(\f,u)^{\a}\big|\lesssim (\sin\f)^{\a}(1-u)^{\max(\a-\b,0)}$} since $f_+(\f,u)=O((1-u))$ by \eqref{f_+rep}. Thus, a~similar argument as the case $\a=0$ shows that $|\pa_m^{\a}H_{\n,+}(m,\f)|\lesssim (\sin\f)^{\a}(1+m\sin\f)^{-\n-\a}$. Using an inequality $\sin\f(1+m\sin\f)^{-1}\leq m^{-1}$, we obtain $(iii)$ for $\a\geq 1$.
\end{proof}

\begin{proof}[Proof of Proposition \ref{Gegenasymp}]
First, we consider the case $\n=0$. From \cite[equation~(4.28)]{BKO}, we have
\begin{align*}
\lim_{\n\to 0}\n^{-1}C_m^{\n}(\cos\f)=\frac{2\cos (m\f)}{m}=\frac{{\rm e}^{{\rm i}m\f}+{\rm e}^{-{\rm i}m\f}}{m}.
\end{align*}
Thus, we can take $g_{\n,\pm}(m,\f)=1/m$ and $r(m,\f)=0$.

Next, we consider the case $\n>0$. For $\f\in \bigl(\frac{\pi}{4}, \frac{3}{4}\pi\bigr)$, we set $g_{\n,\pm}(m,\f)=F_{\n,1}(m)H_{\n,\pm}(m,\f)$ and $r(m,\f)=F_{\n,1}(m)E_{\n}(m,\f)$. For $\f\in \bigl[0,\frac{\pi}{4}\bigr]\cup \bigl[\frac{3}{4}\pi,\pi\bigr]$, we set
\begin{align*}
g_{\n,+}(m,\f)=F_{\n,1}(m)(H_{\n,+}(m,\f)+E_{\n}'(m,\f)),\qquad g_{\n,-}(m,\f)=F_{\n,1}(m)H_{\n,-}(m,\f)
\end{align*}
and $r(m,\f)=0$. Then Proposition \ref{Gegenasymp} directly follows from the identities \eqref{Gegendecom}, \eqref{Gegendecom2}, Proposition \ref{Fnuestimate} and Lemma \ref{Gegenlemma}.
\end{proof}
\begin{thebibliography}{99}
\bibitem[2]{BKO}
Ben~Sa\"{\i}d S., Kobayashi T., {\O}rsted B., Laguerre semigroup and {D}unkl
 operators, \href{https://doi.org/10.1112/S0010437X11007445}{\textit{Compos. Math.}} \textbf{148} (2012), 1265--1336,
 \href{https://arxiv.org/abs/0907.3749}{arXiv:0907.3749}.

\bibitem[17]{IJ}
Ionescu A.D., Jerison D., On the absence of positive eigenvalues of
 {S}chr\"odinger operators with rough potentials, \href{https://doi.org/10.1007/s00039-003-0439-2}{\textit{Geom. Funct. Anal.}}
 \textbf{13} (2003), 1029--1081.

\end{thebibliography}
\end{document}
